\section{国际争端的解决方式}
\begin{enumerate}
  \item 强制方式。
        \begin{enumerate}
          \item 战争或武力。不合法。
          \item 干涉内政。不合法。
          \item 平时封锁。不是绝对合法也不是绝对非法（\zd{安理会做出除外}）。禁止国别封锁。
          \item 反报 vs 报复。针对不违法/违法行为的\zd{对等反措施}。
        \end{enumerate}
  \item 非强制方式。政治性方式、法律解决方式。
        \begin{enumerate}
          \item 政治性方式。\zx{不承担任何法律义务 \hz 后果}。
                \begin{enumerate}
                  \item \zx{翰旋 vs 调停}。第三方帮忙解决争端。提供空间、机会/促使、参与。
                  \item \zx{调查 vs 和解}。第三方\zd{独立解决}争端。解决\zd{事实}/事实 + 法律问题。
                \end{enumerate}
          \item 法律解决方式。仲裁、法院。\zx{单方起诉不受理} \yjt 需要当事方都同意。
                \begin{enumerate}
                  \item 仲裁。国际常设仲裁法院（1900 年，海牙）。影响力变小，了解即可。
                  \item 国际法院。联合国的司法机关。
                        \begin{enumerate}
                          \item 诉讼管辖权。\yi \zd{主权国家有诉权}。\er \zx{判决具有终局性和约束力，但没有强制执行力}（找安理会）。
                          \item 咨询管辖权。\zx{仅对联合国机构开放}，自然人、法人均无权。咨询意见\zd{无}拘束力。
                          \item 其他特殊规定。
                                \begin{itemize}
                                  \item \zx{不适用国籍回避制度，此前参与过是唯一申请回避的理由}。
                                  \item （本国）专案法官制度。
                                  \item \zd{少数服从多数}。
                                \end{itemize}
                        \end{enumerate}
                  \item 国际海洋法法庭。\yi 诉讼主体\zd{扩大到自然人、法人}。\er \zx{对事管辖，受限于海洋争端}。
                \end{enumerate}
        \end{enumerate}
\end{enumerate}
